\section{Structure renders and the geometry of the vacancy site}
\label{si:structures}

The networks read the pristine host supercell of the reference calculation with the oxygen atom to be removed
marked, never the relaxed defect cell. The two renders of this work show that input for two of the screened
oxides, wurtzite ZnO (\artfigoverview{}b of the article) and rutile SnO$_2$
(Figure~\ref{fig:sno2}), and Table~\ref{tab:shell} lists the first cation shell of the removed oxygen in each, read
from the pristine supercells of the release. In wurtzite ZnO the oxygen is tetrahedrally coordinated by four Zn
atoms: three Zn--O bonds of \numRenderZnOFirstShellSeventyFiveDistanceA{}~\AA{} and one of
\numRenderZnOFirstShellZeroDistanceA{}~\AA{}. In rutile SnO$_2$ the oxygen sits in a
planar three-fold site with one Sn--O bond of \numRenderSnOTwoFirstShellZeroDistanceA{}~\AA{} and two of
\numRenderSnOTwoFirstShellEighteenDistanceA{}~\AA{}. Removing the oxygen breaks four and three cation--oxygen bonds respectively, and the
two electrons of the neutral vacancy are left in this cation environment. The
geometry that the network reads is that of the pristine host: in the release's relaxed neutral defect cell of ZnO,
no atom lies more than \numStructureAuditZnoPristineVsRelaxedDisplacementAMaxOfMatchedAtoms{}~\AA{} from its pristine
position, and that relaxation is an output of the defect calculation that the models are meant to replace.

Both renders use the same camera conventions, lights and materials. The SnO$_2$ render was made with Blender
Cycles on the CPU with \numRenderSnOTwoSamples{} samples and denoising, at \numRenderSnOTwoDpi{} dpi over the text
width. Its cell is shifted by lattice translations only, so that the vacancy lies near the faces closest to the
viewer, and the planar three-fold site is seen nearly edge-on; Figure~\ref{fig:sno2} draws the three Sn neighbours
enlarged (parameters in \texttt{figures/si/structure\_sno2\_vacancy.render.json}).

\begin{table}[htbp]
\caption{First cation shell of the removed oxygen in the two rendered hosts (pristine PBEsol(+$U$) supercells of
the release). Atoms: atoms of the supercell; vacancy atom: index of the removed oxygen in the supercell; bond
cutoff: the cation--oxygen distance below which a bond is drawn in the renders. Sources:
\texttt{figures/main/fig1\_zno\_vacancy.render.json} and \texttt{figures/si/structure\_sno2\_vacancy.render.json}
(\texttt{first\_shell}).}
\label{tab:shell}
\centering\small
% Generated by paper/si/tools/make_si_tables.py from figures/main/fig1_zno_vacancy.render.json, figures/si/structure_sno2_vacancy.render.json. Do not edit.
\begin{tabular}{@{}lrrlrr@{}}
\toprule
Host & Atoms & Vacancy atom & Neighbour & M--O distance (\AA) & Bond cutoff (\AA) \\
\midrule
wurtzite ZnO & \numRenderZnONAtoms{} & \numRenderZnOVacancyAtomIndex{} & Zn & \numRenderZnOFirstShellSeventyFiveDistanceA{} & \numRenderZnOCationOBondCutoffAZn{} \\
 &  &  & Zn & \numRenderZnOFirstShellSeventyEightDistanceA{} &  \\
 &  &  & Zn & \numRenderZnOFirstShellOneHundredThirtyFiveDistanceA{} &  \\
 &  &  & Zn & \numRenderZnOFirstShellZeroDistanceA{} &  \\
\addlinespace
rutile SnO$_2$ & \numRenderSnOTwoNAtoms{} & \numRenderSnOTwoVacancyAtomIndex{} & Sn & \numRenderSnOTwoFirstShellZeroDistanceA{} & \numRenderSnOTwoCationOBondCutoffASn{} \\
 &  &  & Sn & \numRenderSnOTwoFirstShellNineteenDistanceA{} &  \\
 &  &  & Sn & \numRenderSnOTwoFirstShellEighteenDistanceA{} &  \\
\bottomrule
\end{tabular}

\end{table}

\begin{figure}[htbp]
\centering
\sifig{structure_sno2_vacancy}
\caption{Pristine supercell of rutile SnO$_2$ from the database (host \mpsnotwo{}, \numRenderSnOTwoNAtoms{} atoms), with the
oxygen to be removed (site O1) and its three nearest Sn neighbours (Table~\ref{tab:shell}) joined into a triangle.
The amber sphere marks the removed oxygen, Sn is grey-blue and O red, and Sn--O bonds shorter than
\numRenderSnOTwoCationOBondCutoffASn{}~\AA{} are drawn except across the cell boundary. Data:
\texttt{figures/si/structure\_sno2\_vacancy.render.json}.}
\label{fig:sno2}
\end{figure}
\FloatBarrier
